% Part: first-order-logic
% Chapter: syntax-and-semantics
% Section: first-order-languages

\documentclass[../../../include/open-logic-section]{subfiles}

\begin{document}

\olfileid{fol}{syn}{fol}

\olsection{प्रथम-क्रम भाषा}


प्रथम-क्रम तर्कशास्त्राची व्यंजके अशा मूलभूत शब्दसंग्रहापासून तयार
केली जातात ज्यात \emph{!!{variable}s}, \emph{!!{constant}s},
\emph{!!{predicate}s} आणि कधीकधी \emph{!!{function}s} असतात. त्यांपासून,
तार्किक संयोजके, संख्यापक आणि कंस व स्वल्पविराम यांसारखी विरामचिन्हे
यांच्या साहाय्याने \emph{पदे} आणि \emph{!!{formula}s} तयार केली जातात.

\begin{explain}
अनौपचारिकपणे, !!{predicate}s ही गुणधर्मांची आणि संबंधांची नावे,
!!{constant}s ही स्वतंत्र वस्तूंची नावे आणि !!{function}s ही
संगतींची नावे आहेत. !!{identity}~$\eq$ सोडल्यास हीच \emph{अतार्किक
चिन्हे} होत आणि ती मिळून एक भाषा बनते. कोणतीही प्रथम-क्रम
भाषा~$\Lang L$ तिच्या अतार्किक चिन्हांनी निर्धारित होते. सर्वांत
सामान्य बाबतीत $\Lang L$ मध्ये प्रत्येक प्रकारची अनंत चिन्हे असतात.
\end{explain}

सामान्य बाबतीत प्रथम-क्रम तर्कशास्त्रात आपण पुढील चिन्हे वापरतो:

\begin{enumerate}
\item तार्किक चिन्हे
\begin{enumerate}
\item तार्किक संयोजके:
  \startycommalist
  \iftag{prvNot}{\ycomma $\lnot$ (नकरण)}{}%
  \iftag{prvAnd}{\ycomma $\land$ (संधियोग)}{}%
  \iftag{prvOr}{\ycomma $\lor$ (विकल्पयोग)}{}%
  \iftag{prvIf}{\ycomma $\lif$ (!!{conditional})}{}%
  \iftag{prvIff}{\ycomma $\liff$ (!!{biconditional})}{}%
  \iftag{prvAll}{\ycomma $\lforall$ (वैश्विक संख्यापक)}{}%
  \iftag{prvEx}{\ycomma $\lexists$ (अस्तित्ववाची संख्यापक)}{}.
\tagitem{prvFalse}{!!{falsity}~$\lfalse$ साठीचे विधानीय स्थिर.}{}
\tagitem{prvTrue}{!!{truth}~$\ltrue$ साठीचे विधानीय स्थिर.}{}
\item द्विस्थानी !!{identity}~$\eq$.
\item !!{variable}s चा एक !!{denumerable}s संच: $\Obj v_0$, $\Obj v_1$, $\Obj
  v_2$, \dots
\end{enumerate}
\item प्रथम-क्रम तर्कशास्त्राची \emph{मानक भाषा} बनवणारी अतार्किक चिन्हे
\begin{enumerate}
\item प्रत्येक $n>0$ साठी $n$-स्थानी !!{predicate}s चा एक !!{denumerable}s संच: $\Obj
  A^n_0$, $\Obj A^n_1$, $\Obj A^n_2$, \dots
\item !!{constant}s चा एक !!{denumerable}s संच: $\Obj c_0$, $\Obj c_1$, $\Obj
  c_2$, \dots.
\item प्रत्येक $n>0$ साठी $n$-स्थानी !!{function}s चा एक !!{denumerable}s संच:
  $\Obj f^n_0$, $\Obj f^n_1$, $\Obj f^n_2$, \dots
\end{enumerate}
\item विरामचिन्हे: (, ) आणि स्वल्पविराम.
\end{enumerate}

आपल्या बहुतेक व्याख्या आणि निष्कर्ष प्रथम-क्रम तर्कशास्त्राच्या संपूर्ण
मानक भाषेसाठी मांडले जातील. मात्र, उपयोगानुसार आपण भाषा केवळ काही
!!{predicate}s, !!{constant}s आणि !!{function}s पुरती मर्यादितही करू शकतो.

\begin{ex}
अंकगणिताची भाषा~$\Lang L_A$ हिच्यात एक द्विस्थानी !!{predicate}~$<$,
एक !!{constant}~$\Obj 0$, एक एकस्थानी !!{function}~$\prime$ आणि दोन
द्विस्थानी !!{function}s~$+$ व~$\times$ असतात.
\end{ex}

\begin{ex}
संचसिद्धान्ताची भाषा~$\Lang L_Z$ हिच्यात केवळ एक द्विस्थानी
!!{predicate}~$\in$ असते.
\end{ex}

\begin{ex}
क्रमांची भाषा~$\Lang L_\le$ हिच्यात केवळ द्विस्थानी
!!{predicate}~$\le$ असते.
\end{ex}

पुन्हा, हे संकेत आहेत: औपचारिकदृष्ट्या ही केवळ पर्यायी नावे आहेत;
उदा., $<$, $\in$ आणि $\le$ ही $\Obj A^2_0$ ची, $\Obj 0$ हे
$\Obj c_0$ चे, $\prime$ हे $\Obj f^1_0$ चे, $+$ हे $\Obj f^2_0$ चे आणि $\times$ हे
$\Obj f^2_1$ चे पर्यायी नाव आहे.

\iftag{defNot,defOr,defAnd,defIf,defIff,defTrue,defFalse,defEx,defAll}{%
वर मांडलेली आदिम संयोजके आणि
\iftag{notprvEx,notprvAll}{संख्यापकाव्यतिरिक्त}{संख्यापकांव्यतिरिक्त}
आपण पुढील \emph{परिभाषित} चिन्हेही वापरतो:
\startycommalist
  \iftag{defNot}{\ycomma $\lnot$ (नकरण)}{}%
  \iftag{defAnd}{\ycomma $\land$ (संधियोग)}{}%
  \iftag{defOr}{\ycomma $\lor$ (विकल्पयोग)}{}%
  \iftag{defIf}{\ycomma $\lif$ (!!{conditional})}{}%
  \iftag{defIff}{\ycomma $\liff$ (!!{biconditional})}{}%
  \iftag{defAll}{\ycomma $\lforall$ (वैश्विक संख्यापक)}{}%
  \iftag{defEx}{\ycomma $\lexists$ (अस्तित्ववाची संख्यापक)}{}%
  \iftag{defFalse}{\ycomma !!{falsity}~$\lfalse$}{}%
  \iftag{defTrue}{\ycomma !!{truth}~$\ltrue$}}{}.

\begin{tagblock}{defNot,defOr,defAnd,defIf,defIff,defTrue,defFalse,defEx,defAll}
\begin{explain}
परिभाषित चिन्ह औपचारिकदृष्ट्या भाषेचा भाग नसते, तर अनौपचारिक संक्षेप
म्हणून मांडले जाते: केवळ आदिम चिन्हे वापरल्यास बरीच लांब होणारी सूत्रे
त्यामुळे संक्षिप्तपणे लिहिता येतात. हा अर्थातच एक लाभ आहे. पण अधिक मोठा
लाभ म्हणजे सिद्धता लहान होतात. एखादे चिन्ह आदिम असल्यास सिद्धतांमध्ये
त्याचा स्वतंत्रपणे विचार करावा लागतो. म्हणून आदिम चिन्हे जितकी अधिक,
तितक्या आपल्या सिद्धता अधिक लांब होतात.
\end{explain}
\end{tagblock}

% Alternate symbols

\begin{intro}
वर आपण वापरलेल्या संज्ञांपेक्षा आणि चिन्हांपेक्षा वेगळ्या संज्ञा व
चिन्हे तुम्हाला परिचित असू शकतात. तर्कशास्त्राच्या ग्रंथांत (आणि
अध्यापनात) ``नकरण'' यासाठी सामान्यतः $\sim$, $\neg$ किंवा~!\ यांपैकी
एखादे, तर ``संधियोग'' यासाठी $\wedge$, $\cdot$ किंवा $\&$ यांपैकी
एखादे चिन्ह वापरतात. ``सशर्त (संकेतार्थक)'' किंवा ``अभिव्यंजन'' यांसाठी
$\rightarrow$, $\Rightarrow$ आणि $\supset$ ही चिन्हे सामान्यतः वापरली जातात.
\iftag{prvIff,defIff}{``द्विसशर्त,'' ``द्वि-अभिव्यंजन'' किंवा
  ``(वास्तविक) सममूल्यता'' यांसाठीची चिन्हे $\leftrightarrow$,
  $\Leftrightarrow$ आणि $\equiv$ ही आहेत.}{}
\iftag{prvFalse,defFalse}{$\lfalse$ या चिन्हाला निरनिराळ्या ठिकाणी
  ``असत्यता,'' ``फाल्सम,'' ``विसंगती'' किंवा ``तळ'' म्हणतात.}{}
\iftag{prvTrue,defTrue}{$\ltrue$ या चिन्हाला निरनिराळ्या ठिकाणी
  ``सत्यता,'' ``वेरम'' किंवा ``शिखर'' म्हणतात.}{}

!!{constant}s (ज्यांना कधीकधी नावे म्हणतात) यांच्यासाठी लॅटिन
वर्णमालेच्या सुरुवातीची लघु अक्षरे (उदा., $a$, $b$, $c$) आणि
!!{variable}s साठी तिच्या शेवटची लघु अक्षरे (उदा., $x$, $y$, $z$)
वापरण्याचा संकेत आहे. संख्यापक !!{variable}s सोबत जोडले जातात, उदा.,
$x$; वैश्विक संख्यापकासाठी $\forall x$, $(\forall x)$, $(x)$, $\Pi x$,
$\bigwedge_x$ आणि अस्तित्ववाची संख्यापकासाठी $\exists x$, $(\exists
x)$, $(Ex)$, $\Sigma x$, $\bigvee_x$ ही संकेतलेखनातील रूपांतरे आहेत.
\end{intro}

\begin{explain}
आपण सर्व विधानीय कारक आणि दोन्ही संख्यापक यांना भाषेची आदिम चिन्हे
मानू शकतो. त्याऐवजी आदिम चिन्हांचा लहान संच निवडून उरलेले
!!{operator}s परिभाषित मानू शकतो. बूलीय कारकांच्या
``सत्यता-फलनात्मकदृष्ट्या संपूर्ण'' संचांत $\{ \lnot, \lor \}$,
$\{ \lnot, \land \}$ आणि $\{ \lnot,
\lif\}$ यांचा समावेश होतो---यांपैकी प्रत्येक संच कोणत्याही एका
संख्यापकासोबत जोडल्यास अभिव्यक्तीदृष्ट्या संपूर्ण प्रथम-क्रम भाषा मिळते.

इतर दोन !!{operator}s तुम्हाला परिचित असू शकतात: शेफर स्ट्रोक~$|$
(हे नाव हेन्री शेफर यांच्यावरून दिले आहे) आणि पिअर्स बाण~$\downarrow$,
ज्याला क्वाइनचा डॅगर असेही म्हणतात. अनुक्रमे ``नँड'' आणि ``नॉर'' हे
त्यांचे नेहमीचे वाचन दिल्यास, हे कारक स्वतःच सत्यता-फलनात्मकदृष्ट्या
संपूर्ण असतात.
\end{explain}

\end{document}
